\section{Limitations}
\label{sec:limitations}
The two-city, three-seed study measures neither training variability nor
statistical significance. Istanbul and New York use different backbone batch
sizes, and full training provenance is unavailable. PO-CFG uses scale $1$
without extrapolation tuning; batch-level distance KL gives no per-sequence
spatial guarantee. Finite-budget projection and decoder fallback leave residual
violations, while same-category refinement cannot repair New York's category-
transition gap. Training-side calibration is not independent validation; its
failed distribution gates and override are reported in \S\ref{sec:exp:tradeoff}
and Appendix~\ref{app:grid}. Stage (d) uses fixed geometry settings, whereas
(d)$\to$(e) jointly changes three factors. Complex order DAGs, conflict
detection, adaptive scheduling, training-time order losses and downstream
recommendation utility remain open. Synthetic output has no formal privacy
guarantee, and distributional similarity does not establish safe release.

These boundaries also delimit the evidence for the control stack. A finite
projection budget can leave requested relations unsatisfied when the backbone
assigns little mass to their categories. The decoder addresses category--POI
mismatch, while category-preserving geometry cannot correct errors already
present in the generated category sequence.

\section{Conclusion}
\label{sec:conclusion}
Order/existence projection, category-consistent decoding, and same-category
refinement form \pcdg on a frozen generator. It outperforms retrained PO-CFG
on all five constraint metrics in both cities and JointGen on six fidelity JSDs
in Istanbul and four in New York. Refinement preserves constraint metrics while
recovering spatial fidelity, but leaves New York's category-transition and
popularity errors unresolved.
The frozen generator can thus serve new category-order requests without
retraining, with spatial refinement restricted to realized categories.
\label{end:main}
